% Draft introduction for the rewritten manuscript (2026-09-20). Plain elsarticle body text; no
% new macros beyond \method{} and \tfe{} from the current preamble.

\section{Introduction}

Multibody systems with lower-pair joints are most conveniently modelled in absolute coordinates:
each body carries its own position and rotation, and every joint contributes algebraic constraints.
The price is an index-3 differential-algebraic system on a Lie group, and the integrator has to keep
three things consistent at once: the constraints at position level, their time derivatives at
velocity and acceleration level, and the rotation group structure of the configuration. Methods
that enforce only the position constraints let the velocity and acceleration constraints drift;
methods that stabilize them by projection or index reduction change the equations that are being
solved; and methods that integrate a reduced set of joint coordinates require the reduced equations
of motion, which are mechanism-specific to derive and lose the direct access to reaction forces
that the absolute formulation provides.

Time finite elements on Lie groups~\cite{chaturvedi2026tfe} have recently shown that high order is
attainable for frictional index-3 pendulum problems, and the absolute-coordinate formulations of
Negrut and co-workers~\cite{taves2020exponential,kissel2021dwelling,kissel2022absolute,fang2023halfimplicit}
have established a public benchmark suite of four mechanisms on which such methods can be compared.
What has been missing is a method of order higher than four for the full index-3 problem in
absolute coordinates whose order can actually be proved rather than observed, and whose proof
survives the details that implementations add: acceleration-level constraint rows, Lagrange
multipliers as stage unknowns, a Newton solver with a finite tolerance, and an endpoint
reconstruction through the exponential map.

This paper studies such a method, \method{}. Its one-step map solves a single square nonlinear
system in which the three-stage Gauss collocation conditions for the joint coordinates, the joint
constraints at position, velocity and acceleration level, and the Newton--Euler balance at every
stage appear as rows (Section~\ref{sec:method}); the rotational endpoint is then reconstructed from
the Gauss quadrature of the body angular velocities through the exponential map. The stage system
has $132$ rows for the two-body chain used throughout, and it is the same system for the four
benchmark mechanisms up to the joint list.

The central result is algebraic. We show that the lifted reduced Gauss stage, that is, the stage
vector obtained by taking the three-stage Gauss collocation solution of the reduced joint-coordinate
equations of motion and lifting it to absolute coordinates, makes every one of the $132$ rows
vanish, and conversely that any root of the non-dynamic rows on the regular branch is such a lift
(Lemma~\ref{lem:exact-stage-identity}). The implemented method is therefore reduced Gauss
collocation, executed in absolute coordinates without ever forming the reduced equations, and the
classical sixth-order theory of Gauss collocation applies to it verbatim. Everything that the
implementation adds on top of the reduced collocation step, namely the endpoint closure of the
velocity-level constraints and the inexact Newton solve, is a perturbation whose size we bound
explicitly. The resulting local error bound (Theorem~\ref{thm:g6fullva-order}) is
$C_{\rm loc}h^7 = (C_G + C_E + C_N c_\eta)h^7$, and its only standing hypothesis is the smoothness
of the constrained problem on a compact neighbourhood of the trajectory: the uniform invertibility
of the stage Jacobian and the reachability of the $c_\eta h^7$ solver tolerance follow from
smoothness for $h\le h_0$ (Lemmas~\ref{lem:p2-from-p1} and~\ref{lem:newton-envelope}). We quantify
$h_0$ on the benchmark and find that it is set by the curvature of the friction law, not by the
collocation structure.

The algebraic identities and the perturbation lemmas are machine-checked in Lean~4 with Mathlib:
the row identities are proved for a transcription of the implemented residual, the perturbation
chain is proved with its explicit constants, and the Gauss tableau is shown to satisfy Butcher's
simplifying assumptions $B(6)$, $C(3)$, $D(3)$ and to fail $B(7)$. Butcher's order theorem and the
classical Gauss local-error estimate enter as cited results. The development ships with the paper
(Appendix~\ref{app:lean}).

Numerically, the method is sixth order on a frictional two-body chain over a full period for six
step sizes (Section~\ref{sec:numerical}), sixth order on the four benchmark mechanisms including
the two closed loops, and its sixth-order cost is repaid in work/precision terms against the
lower-order members of the same Gauss family and against the public absolute-coordinate baseline.
A sweep of the friction regularization locates the regime in which a sharp Stribeck law reduces the
observed order on coarse grids, and the long-time behaviour of the constraint norms and of the
invariants of the frictionless variant is reported. What we do not do is reproduce the temporal
finite-element method of~\cite{chaturvedi2026tfe} on its own problems; the comparison with that
method is at the level of formal order and is stated as such.

\paragraph{Contributions}
\begin{enumerate}
  \item The exact stage identity: the FullVA stage system with acceleration-level constraints is
  reduced three-stage Gauss collocation lifted to absolute coordinates, in both directions.
  \item A sixth-order local and global error theorem for the implemented one-step map under a
  smoothness assumption alone, with the solver and inverse interfaces derived and the small-step
  threshold quantified.
  \item A Lean~4/Mathlib development of $57$ theorems covering the row identities, the perturbation
  chain with explicit constants, and the tableau order conditions.
  \item Sixth-order convergence, work/precision, friction-regime and long-time results on a
  frictional chain and on the four public benchmark mechanisms, all generated from a single
  implementation path and reproducible from the accompanying repository.
\end{enumerate}

Section~\ref{sec:related} reviews Lie-group and index-3 integrators for multibody systems.
Section~\ref{sec:setting} fixes the setting and notation. Section~\ref{sec:method} defines the
one-step map. Section~\ref{sec:order} contains the identity, the perturbation lemmas and the
theorem. Section~\ref{sec:numerical} reports the experiments, Section~\ref{sec:limitations} the
limitations, and Section~\ref{sec:conclusions} concludes.
